\section{Related Work}
\label{sec:related}

\paragraph{Multi-token prediction.} Predicting several future tokens jointly
dates to blockwise-parallel decoding and ProphetNet \citep{qi2020prophetnet},
and was revived as a pretraining objective by \citet{gloeckle2024mtp}, who attach
$n$ \emph{independent output heads} to a shared trunk and report improved sample
efficiency and up to $3\times$ faster inference via self-speculative decoding. Two
properties of their result are central to ours: the heads are \emph{separate}
(each future offset gets its own unembedding), and the benefit is
\emph{increasingly useful for larger model sizes}. DeepSeek-V3
\citep{deepseekv3} adopts sequential MTP modules at frontier scale; Medusa
\citep{cai2024medusa} and related work add prediction heads primarily for
speculative decoding \citep{samragh2025future}; and MuToR \citep{gerontopoulos2025mutor} interleaves
learnable register tokens so a frozen backbone can predict future offsets without
architectural change. Across this line, the auxiliary signal is given \emph{its
own capacity}: separate heads, modules, or register tokens. The present work
studies the opposite, minimal case, a single tied output embedding shared
across horizons, precisely because it isolates the interference that
separate-capacity designs avoid; a separate-head variant would localize whether
the effect is intrinsic to head-sharing (we identify this as the primary future
experiment).

\paragraph{Small-scale degradation of MTP.} \citet{zuhri2025top} document that
multi-token objectives can \emph{hurt} next-token modeling at smaller scales, and
propose token-order prediction as an alternative. This degradation is the failure
mode we instrument; our contribution is not to add to the catalogue of when MTP
helps or hurts, but to measure \emph{what the gradient geometry actually looks
like} at the shared head when it hurts, and to show that the standard conflict
diagnostic mischaracterizes it.

\paragraph{Gradient conflict and surgery.} A large multi-task literature treats
negative inter-task gradient cosine as the signature of interference and
projects it away: PCGrad \citep{yu2020pcgrad}, Gradient Vaccine
\citep{wang2021gradvaccine}, and gradient-similarity loss weighting
\citep{du2018auxsim}, with later refinements \citep{sun2023drmgf,shi2023recon}.
These methods operate on \emph{global} (flattened) gradient vectors and target
gradient \emph{direction}. A parallel line instead balances gradient
\emph{magnitude}: GradNorm \citep{chen2018gradnorm} dynamically tunes per-task
gradient norms rather than reprojecting directions. This magnitude-versus-direction
distinction is central to our findings, which locate the operative structure in
per-row gradient norms rather than in the direction the aggregate cosine reports.
Independently, \citet{jiang2023forkmerge} argue that optimization-based
gradient-coordination methods can overlook auxiliary-target generalization, so
that gradient-conflict signals do not by themselves determine negative transfer;
our measurement gives a concrete mechanism for how a conflict diagnostic can
mislead at the output head. Concretely, our per-row measurement shows the global
cosine is a lossy summary there: it reads $\approx 0$ while the per-row structure
is near-perfectly aligned. Conflict surgery therefore has little of the targeted
directional conflict to act on (Appendix~\ref{app:surgery}), not because surgery
is poorly implemented, but because the aggregate is set by per-row norm structure,
not by a directional conflict surgery can reproject.

\paragraph{The output head as a bottleneck.} \citet{yang2018softmax} showed the
softmax output layer is a rank bottleneck on the distributions a model can
express. \citet{godey2026lmhead} recently showed, for a \emph{single} loss, that
backpropagating through the low-rank output head destroys most of the gradient's
norm, a norm-space single-loss bottleneck. Our result is the direction-space,
\emph{multi-loss} complement: we ask what a cosine diagnostic reports when two
losses share that head, and find their per-row directions align while their
support appears to compete (an inference, not a direct measurement; see
Section~\ref{sec:discussion}).

\paragraph{Positioning.} Relative to prior work, (i) we provide a per-row
gradient measurement at the tied output projection that a global cosine misses;
(ii) we show the next-token degradation of small-scale shared-head multi-token
training, a documented but unexplained failure mode \citep{zuhri2025top}, co-occurs with \emph{aligned} per-row gradients rather than with directional
conflict, which we read as \emph{support} competition (inferred, not directly
measured); and (iii) our capacity-separated auxiliary
\emph{reduces} the next-token cost in the same spirit as capacity-added fixes,
though in our runs it does not fully recover the CE-only baseline within noise
(neither the standard nor a tuned variant); our contribution is not a fix but the
\emph{diagnostic mechanism}: why conflict-based remedies have nothing to act on
here, and why capacity separation is the lever that helps.
